\section{The prepared geometric comparison}
\label{sec:comparison}

The central construction produces a Riemannian exterior whose energy is
almost no larger than the energy of the input, while its full boundary
area is almost no smaller than the input's least enclosing area. This
section states that construction in geometric terms and gives the complete
numerical deduction. The following sections prove the construction. The
data entering it are specified directly; they need not have been obtained
by any particular preparation procedure.

\begin{definition}[Prepared exterior]\label{def:prepared}
Let $(\Omega^n,g_0,K_0)$ be a smooth connected orientable one-ended
exterior with nonempty compact smooth boundary $S$, complete with $S$
included and with compact complement of its Euclidean coordinate end.
Use the constraint and ADM conventions of
Equations~\eqref{eq:setting:constraints}--\eqref{eq:setting:momentum}. Let $\varrho$ be a
smooth positive function on the closed exterior, equal to the coordinate
radius on a distant tail. Require that, for some $0<\beta<1$ and $c>0$,
\begin{equation}\label{eq:end:prepared}
\begin{gathered}
 \operatorname{supp}K_0\text{ is compact},\qquad
 g_0-\delta=O_d(r^{2-n})\text{ for every fixed }d,\\
 R_{g_0}=O(r^{-n-\beta}),\qquad
 \mu_0-|J_0|_{g_0}\ge c\varrho^{-n-\beta},\qquad
 H_S+\operatorname{tr}_S K_0<0\text{ on every component of }S.
\end{gathered}
\end{equation}
The constraint densities are integrable. The ADM momentum is $P_0=0$;
write $E_0$ for the energy. Finally the full-cut infimum, with every
component and all contact counted as in Definition~\ref{def:fullcuts}, is
$A_0=A_{\min,g_0}(S)>0$.
\end{definition}

Once one prepared datum has been fixed, we suppress its subscript and
write $g,K,E,a$ for $g_0,K_0,E_0,A_0$. All constants in its deformation
may depend on this fixed datum. A proper map comparing different prepared
data is an output of end preparation, not an assumption in this definition.

The output below may have a singular frontier. For an enclosing-set
exterior, length distance is computed by
rectifiable paths in the closed exterior and may equal $+\infty$.
Completeness means completeness on each finite-distance equivalence
class; every frontier point and the full perimeter remain included.

\begin{theorem}[The geometric deformation]\label{thm:prepared:deformation}
\label{prop:height:enclosure}
Fix a prepared exterior and $0<\epsilon<1$. For all sufficiently large
integers $N$ there is a connected enclosing-set exterior $X_N$, complete
with its compact frontier $\Sigma_N$ included, in a smooth Riemannian
ambient metric $\widetilde g_N$. It has one Euclidean coordinate end
with compact complement, and
\begin{equation}\label{eq:height:enclosure-estimates}
\begin{split}
 &R_{\widetilde g_N}\ge0,\quad R_{\widetilde g_N}\in L^1(X_N),
 \quad \widetilde g_N-\delta=O_2(r^{2-n}),\\
 &R_{\widetilde g_N}=O(r^{-n-\beta'})\quad\text{for some }\beta'>0,\\
 &E(\widetilde g_N)\le E_0+P_N(e^{-N\epsilon}+e^{-N\epsilon/2}),\\
 &\mathcal H^{n-1}_{\widetilde g_N}(\Sigma_N)
       \ge e^{-(n-1)\epsilon}A_0.
\end{split}
\end{equation}
Here $P_N$ is a polynomial in $N$, with coefficients depending on the
fixed datum and $\epsilon$. The metric is smooth through the frontier.
The whole frontier is outer minimizing and locally perimeter minimizing
as the boundary of its filled side. It is smooth embedded minimal away
from a compact singular set of Hausdorff dimension at most $n-8$; for
$3\le n\le7$ the singular set is empty. Its area counts every component.
\end{theorem}

The proof is completed in Section~\ref{sec:height}. The analytic
construction in Sections~\ref{sec:scalar} and~\ref{sec:elliptic} takes
place on a smooth filling and yields the energy estimate before any
frontier is selected. The height argument then gives a detached minimizing
frontier and the lower bound for its full area.

Here is the precise Riemannian input from the complete companion
\emph{Conformal flow and the Riemannian Penrose inequality with minimizing
frontiers} \cite{RiemannianCompanion}. We use its Theorem~1.1 in the following form.
\begin{quote}
Let $X^n$ be a connected complete enclosing-set exterior, including its
compact frontier $\Sigma$, with a metric $\gamma$ smooth through that
frontier. Suppose there is one Euclidean coordinate end with compact
complement, $\gamma-\delta=O_2(r^{-\tau})$ with
$\tau>(n-2)/2$, and
$R_\gamma\ge0$, $R_\gamma\in L^1$, and
$R_\gamma=O(r^{-n-\beta})$ for some $\beta>0$. Suppose its full frontier
is outer minimizing and locally perimeter minimizing as the boundary
of the filled side, allowing compactly supported variations on either side.
Its regular part is a smooth embedded minimal hypersurface. Its
singular set has Hausdorff dimension at most $n-8$. Then
\[
 E(\gamma)\ge\frac12
 \left(\frac{\mathcal H^{n-1}_\gamma(\Sigma)}{\omega_{n-1}}
 \right)^{(n-2)/(n-1)}.
\]
The assertion includes zero frontier area and imposes neither a spin
condition nor connectedness of the frontier.
\end{quote}

The smooth numerical theorem of Bray--Lee
\cite[Theorem~1.4]{BrayLee2009} is an alternative in dimensions three
through seven. In higher dimensions the companion retains the full
singular-frontier argument. Its conformal-flow ancestry includes Bi--Zhu's
version~2 area-before-nesting and singular separation arguments
\cite{BiZhu2026v2}. The independent smooth positive-mass input is
Brendle--Wang \cite[Corollary~1.6]{BrendleWang2026}, applied there to
constructed smooth complete boundaryless approximants with strictly
positive scalar curvature and the required differentiated end. Those
stronger approximant hypotheses are not additional assumptions on $X$.

\begin{theorem}[The prepared energy inequality]
\label{thm:numerical:prepared}
Every exterior in Definition~\ref{def:prepared} satisfies
\begin{equation}\label{eq:height:prepared-inequality}
 E_0\ge\frac12\left(\frac{A_0}{\omega}\right)^{(n-2)/(n-1)}.
\end{equation}
\end{theorem}
\begin{proof}
Fix the datum and $\epsilon$. Apply the quoted Riemannian theorem to
the actual exterior $X_N$ furnished by
Theorem~\ref{thm:prepared:deformation}. Its metric is smooth through the
compact frontier; its closure is complete and its open exterior is
connected. The end has compact complement and decay exponent
$n-2>(n-2)/2$. Its scalar curvature is nonnegative and integrable with
the required pointwise decay. The whole frontier is outer minimizing
and locally perimeter minimizing, hence has minimal regular part and
singular dimension at most $n-8$. These verify each receiving hypothesis.
We obtain
\[
 E_0+P_N(e^{-N\epsilon}+e^{-N\epsilon/2})
 \ge E(\widetilde g_N)
 \ge\frac12\left(
 \frac{\mathcal H^{n-1}_{\widetilde g_N}(\Sigma_N)}\omega
 \right)^{(n-2)/(n-1)}
 \ge\frac12\left(\frac{e^{-(n-1)\epsilon}A_0}\omega
 \right)^{(n-2)/(n-1)}.
\]
First let $N\to\infty$; exponential decay beats the fixed polynomial.
Then let $\epsilon\downarrow0$. These are limits of numbers. The
exhaustion constructing a smooth solution at each fixed $N$ is completed
before this argument, and no metric at $N=\infty$ is required.
Since $A_0>0$, this also proves $E_0>0$.
\end{proof}

\begin{proof}[Deduction of Theorem~\ref{thm:main}]
Proposition~\ref{prop:end:reduction}, proved in the next section,
supplies prepared data with energies $E_j\to
m=\sqrt{E^2-|P|^2}$ and enclosing infima
$A_j\ge(1-o(1))A_{\min,g}(S)$. Apply
Theorem~\ref{thm:numerical:prepared} to each prepared datum,
completing all its fixed-data limits first. Then
\[
 m\ge\frac12\left(\frac{A_{\min,g}(S)}\omega\right)^{(n-2)/(n-1)}.
\]
No estimate uniform in the preparation index is used.

For sharpness, let $M>0$, $r_M=(2M)^{1/(n-2)}$, and take the static
Schwarzschild--Tangherlini slice
\[
 g=\frac{dr^2}{1-2M/r^{n-2}}+r^2\sigma_{n-1},\quad K=0,
 \qquad r\ge r_M.
\]
Proper distance from the horizon makes this a smooth complete metric
with minimal compact boundary. Its scalar curvature is zero, and its
ADM vector is $(M,0)$. Radial projection to the sphere of radius $r_M$
has $(n-1)$-Jacobian at most one. More explicitly the pullback of that
sphere's area form is closed with comass at most one; Stokes' theorem
on an end truncation of any full outer domain shows that its integral
over the entire cut is $\omega r_M^{n-1}$. Thus every full cut has area
at least that number, while the boundary itself attains it. Consequently
$A_{\min}=\omega(2M)^{(n-1)/(n-2)}$ and equality holds.
\end{proof}

The same numerical bound forces future timelikeness once a positive
enclosing area is retained under an energy-raising conformal change.

\begin{proof}[Proof of Corollary~\ref{cor:automatic-timelikeness}]
If instead $E\le |P|_\delta$, choose $\varepsilon_j>0$ tending to zero
and put
\[
 \lambda_j=\frac{|P|_\delta+\varepsilon_j-E}{2}>0.
\]
Lemma~\ref{lem:end:positive-shell}, proved in the final subsection
of Section~\ref{sec:end}, gives, separately for each $j$,
data satisfying the same one-ended strong-decay, constraint and trapping
hypotheses, with
\[
 E_j=|P|_\delta+\varepsilon_j,\qquad P_j=P,\qquad
 A_{\min,g_j}(S)\ge A_{\min,g}(S)>0.
\]
Their ADM vectors are future timelike, so Theorem~\ref{thm:main} gives
\[
 \sqrt{(|P|_\delta+\varepsilon_j)^2-|P|_\delta^2}
 \ge \frac12\left(\frac{A_{\min,g}(S)}{\omega_{n-1}}\right)^{
                   \frac{n-2}{n-1}}>0.
\]
The right side is fixed and the left side tends to zero, a contradiction.
The theorem is applied afresh to each datum; no estimate for its
end preparation is required to be uniform near the null cone.
Thus $E>|P|_\delta$, and Theorem~\ref{thm:main} applies to the original data.
\end{proof}

The proofs for weak original decay appear in Sections~\ref{sec:weak3}
and~\ref{sec:weak4}. Their reductions reach Definition~\ref{def:prepared}
by different means: three-dimensional enclosing areas converge fully,
whereas the four-dimensional rest replacement gives the one-sided
comparison needed by the numerical argument. The latter is followed by
strictification at fixed replacement radius, with that strictification
removed before the radius tends to infinity.
